\chapter{Databases and data systems}
\label{ch:databases}

\begin{learningobjectives}
After this chapter, you should be able to:
\begin{itemize}
  \item distinguish a database, DBMS, driver, connection, transaction, and schema;
  \item use parameterized queries and explicit transaction boundaries;
  \item compare embedded, remote, warehouse, and object-storage roles; and
  \item choose bounded scans, columnar files, pushdown, or distributed systems
  from measured constraints.
\end{itemize}
\end{learningobjectives}

\begin{chapterconnection}
Arrays and DataFrames usually live inside one process. This chapter crosses a
durability and coordination boundary: data must survive the process, preserve
relationships, support multiple users, and sometimes be queried without being
copied into memory at all.
\end{chapterconnection}

\section{Storage is part of the scientific system}

A \term{database} is organized persisted data plus its logical structure. A
\term{database-management system} (DBMS) is the software that stores, queries,
coordinates, and protects that state. A driver provides a language-specific API
for speaking the DBMS protocol. A connection is one active conversation with
the system.

A \term{schema} names tables, columns, types, constraints, relationships, and
views. It is an executable portion of the data contract. Primary keys express
identity, foreign keys express permitted relationships, uniqueness constraints
prevent selected duplicates, and check constraints reject invalid states.

\section{Transactions preserve invariants}

A transaction groups operations into an atomic success or rollback boundary.
If a workflow inserts an experiment and its measurements, partial success may
create a scientifically misleading state. Transaction ownership should be
visible at the application boundary that knows which operations belong together.

\begin{lstlisting}
with connection:
    connection.execute(
        "INSERT INTO experiments(experiment_id, material) VALUES (?, ?)",
        (experiment_id, material),
    )
    connection.executemany(
        "INSERT INTO measurements(experiment_id, time_s, value) "
        "VALUES (?, ?, ?)",
        rows,
    )
\end{lstlisting}

Values are bound separately from SQL text. String interpolation can change data
into executable syntax, permit injection, and mishandle quoting. Dynamic table
or column identifiers cannot generally use value placeholders; select them from
a small application-owned allowlist.

\section{SQL has semantics, not just syntax}

SQL uses three-valued logic around \code{NULL}. Joins encode relationship
cardinality. Aggregations define grouping and missingness behavior. Window
functions compute across related rows without collapsing them. Query results
require an explicit ordering if order matters.

Before a join, predict whether each input key is unique and how many rows should
result. Afterward, test row counts, key coverage, duplication, and missingness.
A syntactically correct join can create a statistically invalid training table.

\section{Local and remote are deployment shapes}

SQLite is an embedded transactional database stored in a local file. DuckDB is
an in-process analytical database designed for columnar analytics. PostgreSQL
and similar systems run as client/server DBMSs with authentication, concurrent
connections, roles, and network operations. Warehouses and lakehouse engines
serve analytical workloads across managed or distributed storage.

These are not simply size levels of one product. They make different tradeoffs
in concurrency, transactions, operations, latency, governance, and query shape.

\begin{center}
\begin{tikzpicture}[node distance=7mm]
  \node[wideflowbox] (python) {Python analysis\\bounded request};
  \node[wideflowbox,right=of python] (engine) {Query engine\\filter and aggregate};
  \node[wideflowbox,right=of engine] (store) {Database or\\object storage};
  \draw[flowarrow] (python) -- (engine);
  \draw[dataarrow] (store) -- node[above,font=\small]{selected data} (engine);
  \draw[dataarrow] (engine) to[bend left=25] node[below,font=\small]{small result} (python);
\end{tikzpicture}
\end{center}

\section{Object storage is not a filesystem or database}

Cloud object stores organize byte objects by keys and metadata. They are useful
for durable snapshots, partitioned datasets, model artifacts, and large binary
objects. They do not provide ordinary local filesystem semantics or relational
transactions.

An AWS-shaped data platform may combine S3 object storage, a catalog, a query
engine or warehouse, managed relational databases, queues, and identity and
access management. Equivalent roles exist in other clouds and institutional
systems. Learn the role before memorizing a vendor name.

Credentials belong in approved environment configuration, workload identity,
or a secret manager. They must not be embedded in notebooks. Least privilege
means granting the narrow action on the narrow resource for the required time.

\section{When data is too large to load}

Estimate the working set before allocating it. Row count alone is insufficient:
string objects, indexes, temporary arrays, join expansion, and copies can exceed
the apparent file size.

Use the least complex technique that satisfies measured constraints:

\begin{enumerate}
  \item select only necessary rows and columns at the storage/query layer;
  \item use columnar formats such as Parquet for analytical projection;
  \item process bounded batches or iterators;
  \item use lazy plans that push filters and projections toward the source;
  \item use an in-process analytical engine such as DuckDB; and
  \item adopt a distributed system only when one machine cannot meet the
  workload, reliability, or organizational requirement.
\end{enumerate}

Distributed computation adds serialization, partitioning, shuffle, scheduling,
failure recovery, observability, and cost. It is an architecture, not a badge of
seriousness.

\begin{checkpointbox}
A 60 GB Parquet dataset contains 120 columns, but an analysis needs six columns
and one month of records. Propose the first three approaches you would try,
state what evidence each collects, and explain when a cluster becomes justified.
\end{checkpointbox}

\section{Worked example: from a measurement contract to SQL}

Suppose a laboratory records experiments and temperature measurements. Begin
with meaning rather than table syntax:

\begin{itemize}
  \item one experiment has one stable identifier and one material;
  \item each measurement belongs to exactly one experiment;
  \item time is measured in seconds from the experiment start;
  \item temperature is stored in degrees Celsius; and
  \item an experiment cannot contain two measurements at the same time.
\end{itemize}

The schema can enforce part of this contract.

\begin{workedexample}{translate scientific rules into constraints}
\begin{lstlisting}[language=SQL]
CREATE TABLE experiments (
    experiment_id TEXT PRIMARY KEY,
    material TEXT NOT NULL
);

CREATE TABLE measurements (
    experiment_id TEXT NOT NULL,
    time_s REAL NOT NULL CHECK (time_s >= 0),
    temperature_c REAL NOT NULL,
    PRIMARY KEY (experiment_id, time_s),
    FOREIGN KEY (experiment_id)
        REFERENCES experiments(experiment_id)
);
\end{lstlisting}

\textbf{Step 1: identify rows.} One row in \code{experiments} is one
experiment. One row in \code{measurements} is one reading at one experiment
time.

\textbf{Step 2: identify identity.} The composite primary key says that the
pair \code{(experiment\_id, time\_s)} uniquely identifies a measurement.

\textbf{Step 3: identify valid relationships.} The foreign key prevents a
measurement from referring to an experiment that does not exist.

\textbf{Step 4: identify what remains outside SQL.} The database does not know
whether the sensor was calibrated, whether Celsius is the correct source unit,
or whether the experiment clock was reset. Application validation and
scientific review still have work to do.
\end{workedexample}

Now consider the question, ``What is the mean temperature for each material?''
The query must join the tables because material and temperature live in
different relations.

\begin{lstlisting}[language=SQL]
SELECT e.material, AVG(m.temperature_c) AS mean_temperature_c
FROM experiments AS e
JOIN measurements AS m
  ON m.experiment_id = e.experiment_id
GROUP BY e.material
ORDER BY e.material;
\end{lstlisting}

Before running it, predict the join cardinality: each measurement should match
exactly one experiment. After running it, compare the joined row count with the
measurement row count. This does not prove the scientific interpretation, but
it tests an important structural assumption.

\section{Missingness and joins deserve explicit examples}

SQL \code{NULL} means missing or unknown at the database level; it is not zero,
an empty string, or the floating-point value \code{NaN}. The comparison
\code{value = NULL} does not produce true or false in the ordinary way. Use
\code{value IS NULL}. Aggregates also differ: \code{COUNT(*)} counts rows,
whereas \code{COUNT(value)} counts non-NULL values.

A left join keeps every row from the left relation. If no right-side match
exists, right-side columns become \code{NULL}. An inner join silently removes
unmatched left rows. Neither choice is inherently correct; the research
question determines whether absent matches should disappear, remain missing,
or trigger an error.

\begin{misconceptionbox}
\textbf{Misconception:} a database is a large CSV file with faster search.

\textbf{Correction:} a DBMS coordinates typed structure, constraints,
transactions, concurrent users, recovery, permissions, and query planning. A
CSV file stores bytes. Both can be useful, but they provide different
guarantees and failure modes.
\end{misconceptionbox}

\section{A measured decision process for large data}

``Does not fit in memory'' is a diagnosis to quantify, not a command to adopt a
cluster. Record available memory, compressed file size, projected columns,
selected rows, intermediate operations, and acceptable runtime. Then test the
smallest architecture that could work.

For the 60 GB checkpoint dataset, first inspect metadata without loading all
values. Next ask a query engine to project six columns and push the date filter
into the scan. Finally measure the result and its largest intermediate. If one
machine can execute that plan reliably, distribution adds coordination cost
without solving a demonstrated problem. If the filtered working set, required
join, concurrency, or deadline exceeds one machine, the measurement explains
why a larger system is justified.

\conceptcheck{Why can a 4 GB CSV require substantially more than 4 GB of memory
as a pandas DataFrame? Name at least three sources of additional memory use.}

\section{Exercises}

\subsection*{Foundational}
\begin{enumerate}
  \item Distinguish a database, DBMS, schema, driver, connection, and
  transaction using the laboratory example.
  \item Explain why query result order is unspecified without
  \code{ORDER BY}.
  \item Compare \code{COUNT(*)}, \code{COUNT(temperature\_c)}, and
  \code{AVG(temperature\_c)} when some temperatures are NULL.
\end{enumerate}

\subsection*{Applied}
\begin{enumerate}
  \item Add a sensor table and a foreign key from measurements. State the row
  meaning and identity rule for each table before writing SQL.
  \item Construct a join that accidentally multiplies rows. Write precondition
  and postcondition queries that reveal the error.
  \item Design a parameterized query that returns a bounded time interval and
  only two analytical columns. Explain which work should occur in the DBMS.
\end{enumerate}

\subsection*{Design}
Choose a laboratory or industry dataset. Propose a local prototype architecture
and a remote production architecture. For each, state the storage role,
transaction needs, access controls, expected working set, failure recovery, and
the measurement that would justify moving to a distributed system.

\begin{chapterreview}
\textbf{Key ideas.} A schema is executable data meaning. Transactions preserve
multi-step invariants. Parameter binding separates values from executable SQL.
Joins and missingness can change a population without producing an error.
Scalable analysis begins with projection, filtering, and measured working sets.

\textbf{Vocabulary.} database; DBMS; schema; primary key; foreign key;
constraint; transaction; rollback; parameterized query; NULL; join cardinality;
query pushdown; columnar storage; object storage; distributed system.

\textbf{Explain aloud.} Why can a query be syntactically valid, complete
successfully, and still create an invalid machine-learning training table?
\end{chapterreview}

\begin{notebooklink}
Lecture 3 notebook 05 builds a constrained SQLite system, connects pandas and
SQLAlchemy, queries with DuckDB, scans Arrow and Polars data lazily, and maps the
same roles onto AWS-style workflows without requiring cloud credentials.
\end{notebooklink}

\section{Takeaway}

Data systems make identity, relationships, durability, concurrency, and access
operational. Correct queries begin with row meaning and constraints. Scalable
work begins by reducing unnecessary data movement.
